\podnaslov{Ocenjevanje in napovedovanje}

Naše opažanje sestoji iz nekih podatkov $X$.
Zanima nas druga količina $Y$, ki bi jo radi vsaj približno ocenili.
Postavimo \pojem{matematični model}, ki je opis mehanizma, za katerega
domnevamo, da je generiral $X$ in $Y$.
V statistiki vzamemo verjetnostni model, kjer je vsaj ena od $X$ oz.~$Y$
slučajna.
Iščemo opazljivo količino $\hat{Y}$ (tj.~tako, ki jo lahko izračunamo z merljivo
funkcijo $h(X)$), ki bo čim bližje $Y$.

V Bayesovem modelu sta $X$ in $Y$ slučajni spremenljivki na fiksnem
verjetnostnem prostoru $(\Omega, \mathcal{F}, P)$.
Če poznamo $X$ in $Y$ živi na množici $\R$, je smiselna napoved $\hat{Y} = E(Y
\such X)$.
Postopek je naslednji:
\begin{itemize}
\item Poznamo $P(X = x \such Y = y)$
\item Privzamemo apriorno porazdelitev $Y$
\item Po Bayesovi formuli izračunamo $P(Y = y \such X = x)$
\item Izračunamo $E(Y \such X = x)$
\end{itemize}
Bayesova statistika je pripravna, če podatke dobivamo postopoma.

Drug koncept je postavil Ronald Fisher.
Tu imamo več verjetnostnih mer $P_\theta$ za $\theta \in I$.
Modeliramo tako, da je $X$ slučajna spremenljivka, ki jo poznamo, $y =
g(\theta)$ pa je deterministična in je ne poznamo.
Karakteristiko $y$ ocenimo z njeno \pojem{cenilko} $\hat{y}$.
O cenilki govorimo, preden opazimo podatke; ko jih poznamo, iz njih izračunamo
vrednost cenilke, in ji pravimo \pojem{ocena}.

\vprasanje{Kakšna je razlika med Bayesovo in Fisherjevo statistiko?}

\podnaslov{Vrednotenje cenilk in predikatorjev}

Naj bo $\hat{y} = h(X)$ cenilka za $y$.
\pojem{Pričakovana} ali \pojem{srednja kvadratična napaka} cenilke je
\[
  \operatorname{MSE}_\theta(\hat{y} \such y) = E_\theta((\hat{y} - y)^2).
\]
Pogosto jo še korenimo, da dobimo RMSE.
Definiramo tudi \pojem{pristranskost}
\[
  \operatorname{Bias}_\theta(\hat{y} \such y)
  = E_\theta(\hat{y}) - y.
\]
Cenilka je \pojem{nepristranska}, če je pristranskost enaka $0$ za vse $\theta$.
Opazimo, da je za nepristransko cenilko pričakovana kvadratna napaka enaka
varianci.
V tem primeru količini RMSE pravimo \pojem{standardna napaka}.
V splošnem je
\[
  \var_\theta(\hat{y})
  = \var_\theta(\hat{y} - y)
  = E_\theta((\hat{y} - y)^2) - (E_\theta(\hat{y} - y))^2
  = \operatorname{MSE}_\theta(\hat{y} \such y)
  - \left( \operatorname{Bias}(\hat{y} \such y) \right)^2.
\]

\vprasanje{Definiraj pričakovano napako in pristranskost.}

Če podatke dobivamo postopoma, dobimo zaporedje cenilk $\hat{y}_n$.
Zaporedje je \pojem{šibko dosledno}, če gre
\[
  \hat{y}_n \xrightarrow[n \to \infty]{d \such \theta} y
\]
za vse $\theta \in I$.

\begin{trditev}
  Če je
  \[
	\lim_{n \to \infty} \operatorname{MSE}_\theta(\hat{y}_n \such y) = 0
  \]
  za vse $\theta \in I$, je zaporedje $(\hat{y}_n)_n$ šibko dosledno.
\end{trditev}

\begin{proof}
  Sledi iz ene od trditev od prej; če je
  \[
	\lim_{n \to \infty} E((d(X_n, c))^p) = 0,
  \]
  potem $X_n \xrightarrow[n \to \infty]{d} c$.
\end{proof}

\vprasanje{Kdaj je zaporedje cenilk šibko dosledno? Povej zadostni pogoj.}

\begin{posledica}
  Če je zaporedje cenilk $(\hat{y}_n)_n$ asimptotično nepristransko, tj.~da
  velja $\lim E_\theta(\hat{y}_n) = y$ za vsak $\theta \in I$, in če je $\lim
  \var_\theta(\hat{y}_n) = 0$, je to zaporedje šibko dosledno.
\end{posledica}

Če so $X_1, X_2, \ldots$ enako porazdeljene z $E(X_1) = \mu$ in $E(X_1^2) <
\infty$ ter nekorelirane, je $\overline{X}$ dosledna cenilka za $\mu$.
Poleg tega je $\overline{X}$ \pojem{najboljša nepristranska linearna cenilka}
(NNLC) za $\mu$.
Najboljša tu pomeni, da ima najmanjšo srednjo kvadratno napako.
Pri tem pazimo, da je $\overline{X}^2$ pristranska cenilka za $\mu^2$; velja
\[
  E(\overline{X}^2) = \var(\overline{X}) + (E(\overline{X}))^2
  = \frac{\sigma^2}{n} + \mu^2,
\]
je pa asimptotično nepristranska in šibko dosledna.

% LocalWords:  generiral Bayesovem živi Bayesovi Bayesova Fisher Fisherjevo
% LocalWords:  Bayesovo predikatorjev korenimo RMSE NNLC nekorelirane
